\documentclass[10pt,twocolumn]{article}
\usepackage[margin=0.85in]{geometry}
\IfFileExists{lmodern.sty}{\usepackage[T1]{fontenc}\usepackage{lmodern}\usepackage{microtype}}{}
\usepackage{amsmath}
\usepackage{booktabs}
\usepackage{xcolor}
\usepackage{pgfplots}
\pgfplotsset{compat=1.17}
\usepackage{hyperref}
\usepackage{fancyhdr}
\hypersetup{colorlinks=true,linkcolor=black,citecolor=black,urlcolor=black,
  pdftitle={Does Adaptive Checkpointing Help Spark Structured Streaming? A Fault-Injection Study of Replay, Recovery and Latency},
  pdfauthor={Rajshekar Medipally},
  pdfkeywords={stream processing, fault tolerance, checkpointing, Apache Spark, Structured Streaming, fault injection, negative result}}

\pagestyle{fancy}
\fancyhf{}
\fancyfoot[C]{\thepage}
\renewcommand{\headrulewidth}{0pt}

% Every number in this paper is generated by paper/make_tables.py from the trial
% records committed under experiments/; in-text numbers come from tables/macros.tex.
% generated by paper/make_tables.py from the committed trial records; do not edit
\newcommand{\SOneTrials}{84}
\newcommand{\SOneCommits}{4a2dfe3, 5607d08}
\newcommand{\SOneRecMin}{9.5}
\newcommand{\SOneRecMax}{28.6}
\newcommand{\SOneCorruptFailed}{40}
\newcommand{\SOneCorrupt}{40}
\newcommand{\SOneStartupShareMed}{67}
\newcommand{\SOneStartupShareMin}{57}
\newcommand{\SOneStartupShareN}{40}
\newcommand{\SimCorruptDupB}{2.96}
\newcommand{\SimCorruptDupA}{0.18}
\newcommand{\SimAcsDup}{0.066}
\newcommand{\PairPfiftyAcsTwoOne}{6/10}
\newcommand{\PairPninetyfiveAcsTwoOne}{3/10}
\newcommand{\PairCommitsAcsTwoOne}{0/10}
\newcommand{\PairPfiftyAcsOneOne}{9/10}
\newcommand{\PairPninetyfiveAcsOneOne}{7/10}
\newcommand{\PairCommitsAcsOneOne}{2/10}
\newcommand{\PairPfiftyAcsTwoTwo}{7/10}
\newcommand{\PairPninetyfiveAcsTwoTwo}{1/10}
\newcommand{\PairCommitsAcsTwoTwo}{1/10}
\newcommand{\PairPfiftyAcsOneTwo}{5/10}
\newcommand{\PairPninetyfiveAcsOneTwo}{0/10}
\newcommand{\PairCommitsAcsOneTwo}{2/10}
\newcommand{\DiffPfiftyAcsTwoOne}{$-$0.27 {[}$-$0.52, $-$0.02{]}}
\newcommand{\DiffPninetyfiveAcsTwoOne}{+0.90 {[}$-$0.58, +2.37{]}}
\newcommand{\DiffCommitsAcsTwoOne}{+8.2 {[}+4.9, +11.5{]}}
\newcommand{\DiffPfiftyAcsOneOne}{$-$0.68 {[}$-$1.35, $-$0.01{]}}
\newcommand{\DiffPninetyfiveAcsOneOne}{$-$1.48 {[}$-$3.13, +0.17{]}}
\newcommand{\DiffCommitsAcsOneOne}{+6.9 {[}+2.2, +11.6{]}}
\newcommand{\DiffPfiftyAcsTwoTwo}{$-$0.02 {[}$-$0.47, +0.44{]}}
\newcommand{\DiffPninetyfiveAcsTwoTwo}{+1.14 {[}+0.09, +2.20{]}}
\newcommand{\DiffCommitsAcsTwoTwo}{+6.7 {[}+3.0, +10.4{]}}
\newcommand{\DiffPfiftyAcsOneTwo}{+0.04 {[}$-$0.80, +0.88{]}}
\newcommand{\DiffPninetyfiveAcsOneTwo}{+2.71 {[}+1.43, +3.99{]}}
\newcommand{\DiffCommitsAcsOneTwo}{+5.2 {[}+1.4, +9.0{]}}
\newcommand{\FirstOutMed}{11}
\newcommand{\FirstOutMin}{9.5}
\newcommand{\FirstOutMax}{21.6}
\newcommand{\KillSpreadMed}{2.3}
\newcommand{\KillSpreadMax}{12.1}
\newcommand{\CurveAcsOneOne}{2.56}
\newcommand{\CurveAcsTwoOne}{1.93}
\newcommand{\CurveAcsOneTwo}{2.73}
\newcommand{\CurveAcsTwoTwo}{2.02}
\newcommand{\ACSStartupShareMed}{67}
\newcommand{\FastStartup}{6.7}
\newcommand{\FastFirst}{3.3}
\newcommand{\SlowStartup}{9.1}
\newcommand{\SlowFirst}{5.0}
\newcommand{\SwitchKilled}{2}
\newcommand{\SwitchKilledMax}{0.1}
\newcommand{\ACSTrialsOne}{60}
\newcommand{\ACSTrialsTwo}{66}
\newcommand{\ACSCleanOne}{60}
\newcommand{\ACSCleanTwo}{60}
\newcommand{\ACSStalled}{6}
\newcommand{\ACSCommits}{369028b, 3d1f93f, 7d66520}
\newcommand{\DirtyTrials}{0}
\newcommand{\AllTrials}{210}
\newcommand{\OverrunUnderOne}{120}
\newcommand{\OverrunTotal}{126}
\newcommand{\ReplayEvents}{9}
\newcommand{\CleanTrials}{120}
\newcommand{\DupTrials}{1}
\newcommand{\DupRows}{1{,}739}
\newcommand{\ACSTrialsClean}{40}
\newcommand{\Switches}{320}
\newcommand{\SwitchN}{318}
\newcommand{\SwitchMed}{0.54}
\newcommand{\SwitchMin}{0.52}
\newcommand{\SwitchMax}{2.07}
\newcommand{\SwitchOverOne}{33}
\newcommand{\RecFast}{90}
\newcommand{\RecSlow}{29}
\newcommand{\RecOutlier}{1}
\newcommand{\RecAll}{120}
\newcommand{\ACSDupTrials}{0}
\newcommand{\PfiftyAOne}{0.76}
\newcommand{\PninetyfiveAOne}{6.10}
\newcommand{\CommitsAOne}{84}
\newcommand{\PfiftyCthreeOne}{2.12}
\newcommand{\PninetyfiveCthreeOne}{7.01}
\newcommand{\CommitsCthreeOne}{30}
\newcommand{\PfiftyAcsTwoOne}{1.73}
\newcommand{\PninetyfiveAcsTwoOne}{8.36}
\newcommand{\CommitsAcsTwoOne}{37}
\newcommand{\IntervalAcsTwoOne}{3.3}
\newcommand{\PfiftyCfiveOne}{3.96}
\newcommand{\PninetyfiveCfiveOne}{11.92}
\newcommand{\CommitsCfiveOne}{17}
\newcommand{\PfiftyAcsOneOne}{3.48}
\newcommand{\PninetyfiveAcsOneOne}{10.27}
\newcommand{\CommitsAcsOneOne}{26}
\newcommand{\IntervalAcsOneOne}{5.8}
\newcommand{\PfiftyBOne}{12.65}
\newcommand{\PninetyfiveBOne}{29.23}
\newcommand{\CommitsBOne}{6}
\newcommand{\PfiftyATwo}{0.84}
\newcommand{\PninetyfiveATwo}{8.75}
\newcommand{\CommitsATwo}{82}
\newcommand{\PfiftyCthreeTwo}{2.18}
\newcommand{\PninetyfiveCthreeTwo}{8.61}
\newcommand{\CommitsCthreeTwo}{28}
\newcommand{\PfiftyAcsTwoTwo}{1.99}
\newcommand{\PninetyfiveAcsTwoTwo}{8.89}
\newcommand{\CommitsAcsTwoTwo}{35}
\newcommand{\IntervalAcsTwoTwo}{3.5}
\newcommand{\PfiftyCfiveTwo}{3.46}
\newcommand{\PninetyfiveCfiveTwo}{8.11}
\newcommand{\CommitsCfiveTwo}{18}
\newcommand{\PfiftyAcsOneTwo}{3.68}
\newcommand{\PninetyfiveAcsOneTwo}{10.88}
\newcommand{\CommitsAcsOneTwo}{24}
\newcommand{\IntervalAcsOneTwo}{5.8}
\newcommand{\PfiftyBTwo}{13.98}
\newcommand{\PninetyfiveBTwo}{30.44}
\newcommand{\CommitsBTwo}{6}

\makeatletter
\def\StudyOneRows{\@@input tables/study1.tex }
\def\SimRows{\@@input tables/sim_vs_measured.tex }
\def\RunTwoRows{\@@input tables/acs_run2.tex }
\def\BothRows{\@@input tables/acs_both.tex }
\def\PairRows{\@@input tables/paired.tex }
\def\ReplayRows{\@@input tables/replay.tex }
\def\SplitRows{\@@input tables/recovery_split.tex }
\def\StallRows{\@@input tables/stalls.tex }
\makeatother

\title{Does Adaptive Checkpointing Help Spark Structured Streaming?\\
\large A Fault-Injection Study of Replay, Recovery and Latency}
\author{Rajshekar Medipally\\
\small Independent researcher\\
\small \url{https://github.com/rmedipallycic/spark-streaming-fault-tolerance}}
\date{October 2026}

\begin{document}
\maketitle
\thispagestyle{fancy}

\begin{abstract}
Checkpoint-interval tuning is usually framed as a trade-off between overhead
in normal operation and work lost on failure, and adaptive controllers aim to
shorten the interval when failures look likely. We tested whether that framing
holds for Spark Structured Streaming, where the micro-batch trigger interval
also sets how often the offset and commit logs are written. We built a
fault-injection harness that SIGKILLs a Spark job reading sequence-numbered
records from Kafka, restarts it from its checkpoint, and counts every
duplicated or missing record, and implemented an adaptive controller (ACS)
that changes the trigger at runtime by restarting the query between batches.
Over \AllTrials{} measured trials on one machine, a crash caused a batch to be
reprocessed in only \ReplayEvents{} of \CleanTrials{} comparison trials, because
records that arrive between batches are read for the first time after a
restart, not replayed; recovery took about 10\,s whatever the interval and was
dominated by restart. The interval therefore traded end-to-end latency against
checkpoint overhead, not against lost work. Interval switches were cheap (median
\SwitchMed\,s) and safe (no duplicates in \Switches{} switches), but on paired,
identical workloads and failures a fixed interval chosen at ACS's own mean
matched or beat it: against a fixed 5.8\,s interval, the original ACS rule had
the lower 95th-percentile latency in \PairPninetyfiveAcsOneTwo{} pairs and
fewer commits in only \PairCommitsAcsOneTwo. The original rule also cannot react
to any single signal. We report the negative result, its mechanism, and two
measurement pitfalls we hit along the way: an earlier simulation that predicted
the wrong failure behaviour, and laptop standby that silently paused six
trials, which we detected from the producer's clock.
\end{abstract}

\section{Introduction}

Fault-tolerant stream processors periodically record how far they have read and
what state they hold, so that after a failure they can resume instead of
starting over. How often to do this is a classic question
\cite{young1974first,daly2006higher}: checkpoint too often and normal operation
pays for it; too rarely and a failure throws away more work. Adaptive policies
that tighten the interval when the system looks unstable follow naturally from
that framing.

We set out to build and evaluate such a policy for Spark Structured Streaming
\cite{armbrust2018structured}. An earlier simulation of the policy, Adaptive
Checkpoint Selection (ACS), suggested large reductions in duplicated records.
When we audited that simulation, we found that its advantage followed from how
the model was written, and we replaced it with live measurements on real Spark.
This paper reports what those measurements showed.

The answer was negative, for a reason that is specific to micro-batch execution
and easy to miss. In Structured Streaming the processing-time trigger sets both
how often a batch runs and how often progress is written to the checkpoint.
Records that arrive between two triggers are not lost or replayed when the
driver dies: they are simply read for the first time by the restarted query.
Only a batch that is \emph{running} at the moment of the crash is redone. Longer
intervals therefore do not mean more lost work in the usual sense; they mean
higher latency, because records wait longer for the next batch. The interval is
a latency-versus-overhead knob, and a well-chosen fixed value served that
trade-off as well as adaptive control did.

\paragraph{Contributions.}
\begin{enumerate}
\item A fault-injection harness for Spark Structured Streaming with Kafka that
counts every duplicated and missing record by sequence number, measures replay
from the checkpoint logs, end-to-end latency per record, and recovery time, and
records the code commit and environment of every trial.
\item A measured account of delivery under driver failure and checkpoint
corruption for the exactly-once file sink and an at-least-once foreachBatch
sink (\SOneTrials{} trials).
\item A live, restart-based implementation of adaptive trigger control, with
its switching cost, and an analysis showing that the original ACS rule cannot
react to any single signal.
\item A paired comparison of ACS against fixed intervals, including fixed
intervals set to ACS's own mean, in two runs (\ACSTrialsOne{} and
\ACSTrialsTwo{} trials), showing no advantage for adaptation and explaining why.
\item Two methodology lessons: simulated failure behaviour that differed from
the measured behaviour in kind, and a clock-based test that detected trials
paused by laptop standby.
\end{enumerate}

\section{Background}

\paragraph{Micro-batch execution.}
Structured Streaming runs a streaming query as a sequence of small batch jobs
\cite{armbrust2018structured,zaharia2013discretized}. With a processing-time
trigger of $T$ seconds, the driver starts a batch every $T$ seconds if new data
is available. Before a batch runs, the driver writes the range of source
offsets it will read to the checkpoint's \texttt{offsets/} log; after the sink
has finished, it writes an entry to \texttt{commits/}. On restart, Spark finds
the latest offsets entry; if it has no matching commit, the batch is run again
with the same offsets.

\paragraph{End-to-end guarantees.}
Exactly-once output needs a replayable source (Kafka offsets
\cite{kreps2011kafka}), deterministic re-execution of a batch from its recorded
offsets, and an idempotent or transactional sink \cite{sparkguide}. Spark's file
sink provides the last part through its own \texttt{\_spark\_metadata} log,
which lists the files that belong to each committed batch. A
\texttt{foreachBatch} sink hands each batch to user code; if that code is not
idempotent, a batch that is run again is written again (at-least-once).

\paragraph{What the interval controls.}
Because offsets and commits are written once per batch, the trigger interval is
also the checkpoint interval. Section~\ref{sec:mechanism} shows why, unlike in
the classic model, a longer interval does not proportionally increase the work
redone after a crash.

\section{Harness and setup}
\label{sec:harness}

\paragraph{One trial.}
A trial creates a fresh Kafka topic and produces 36{,}000 records with
sequence numbers through an idempotent producer, so that any duplicate is
Spark's. It starts the Spark job from the earliest offset, waits for the first
committed output, injects the failure 15\,s plus a per-trial phase later,
restarts the job from the same checkpoint, waits until output stops changing,
and stops the query gracefully. It then counts duplicated and missing sequence
numbers in the committed output and, separately, in every output file. Each
trial writes a JSON record with its configuration, git commit (and whether the
working tree was dirty; \DirtyTrials{} trials were run from a dirty tree),
package versions, event timestamps and counts, plus the Spark log, the
controller log and a checkpoint snapshot taken before the restart.

\paragraph{Failures.}
\emph{Driver kill}: SIGKILL of the Spark process group (in local mode, driver
and executors are one JVM), then restart. \emph{Offset-log corruption}: the same
kill, then random bytes written over the middle of the latest file in
\texttt{offsets/}, then restart.

\paragraph{Metrics.}
\emph{Recovery}: time from the kill to the first new entry in \texttt{commits/}
after the restart. \emph{Replay}: records in batches that had an offsets entry
but no commit at the kill, read from the checkpoint before the restart.
\emph{End-to-end latency}: per record, the modification time of the first output
file containing it minus the producer's send time; we report each trial's
median (p50) and 95th percentile (p95). \emph{Batch commits}: the number of
committed micro-batches in the trial, our measure of checkpoint overhead (every
trial processes the same 36{,}000 records).

\paragraph{Environment.}
One laptop (Intel Core i5-1335U) running WSL2 Ubuntu with 8\,GB for WSL;
Spark Structured Streaming 3.5.3 through PySpark in \texttt{local[2]} mode;
Java 17; one Kafka broker (Confluent 7.5.0) in Docker Compose, with a
4-partition topic per trial. Everything ran on the same machine, so network and
cluster scheduling play no part; Section~\ref{sec:threats} discusses what this
does and does not allow us to conclude.

\section{Study 1: what Spark delivers under failure}
\label{sec:study1}

The first study used two fixed triggers, A (1\,s) and B (30\,s), both sinks,
both failures, and 10 trials per cell plus one failure-free baseline each, at a
constant 300 records/s (\SOneTrials{} trials, code commits \SOneCommits; an
earlier attempt at B, invalidated by a full \texttt{/tmp}, is kept and
documented in the repository). Kill times were spread across the trigger
interval by a golden-ratio sequence.

\begin{table*}[t]
\centering\small
\caption{Study 1: driver kill and offset-log corruption (10 trials per cell;
baselines: 1). Missing records are counted only for trials whose restart
succeeded. Recovery: median (min--max), seconds.}
\label{tab:study1}
\resizebox{\textwidth}{!}{%
\begin{tabular}{@{}lllrlrrrl@{}}
\toprule
Failure & Sink & Trigger & Trials & Outcomes & Trials w/ dup. & Max dup. rows & Missing & Recovery s \\
\midrule
\StudyOneRows
\bottomrule
\end{tabular}}
\end{table*}

Table~\ref{tab:study1} shows three results. The file sink delivered every
record exactly once in all 20 driver kills; even the uncommitted files contained
no duplicates. The at-least-once foreachBatch sink duplicated records in one
trial: 300 records, exactly one 1-second batch at 300 records/s, written again
after the restart. A duplicate requires the kill to land between a batch's write
and its commit, a short window that none of the ten 30\,s kills hit. Offset-log
corruption stopped the pipeline in \SOneCorruptFailed{} of \SOneCorrupt{}
trials: Spark refused to restart (\texttt{MalformedInputException}, exit code
3) instead of continuing with wrong offsets, so no records were duplicated or
silently lost; processing simply stopped until an operator intervened.

Recovery after a driver kill took \SOneRecMin--\SOneRecMax\,s. From the
timestamps in each Spark log, a median \SOneStartupShareMed\% (at least
\SOneStartupShareMin\%, over \SOneStartupShareN{} kills) of that time passed
before the restarted query was running, i.e.\ in JVM and Spark start-up.

\paragraph{Simulation versus measurement.}
Table~\ref{tab:sim} compares these measurements with the earlier simulation of
the same strategies. The simulation assigned every failure a constant duplicate
rate, including \SimCorruptDupA\% (A) and \SimCorruptDupB\% (B) for checkpoint
corruption, where the measured pipeline instead stopped. It had no model of a
sink, so it could not represent the file sink, which admits no duplicates by
design, and it predicted recovery growing about fivefold from A to B. Measured, B
recovered a few seconds slower than A in this study, but A and B ran on
different days; when every strategy was interleaved (Section~\ref{sec:study2}),
recovery did not depend on the interval. The differences are in kind, not degree, and they are
why this paper reports only measured results.

\begin{table}[t]
\centering\small
\caption{Driver failure: earlier simulation versus Study 1. Duplicate rate:
duplicated rows as a share of records, mean over trials; fBatch: foreachBatch sink. Recovery: simulated
mean versus measured median, both sinks.}
\label{tab:sim}
\setlength{\tabcolsep}{3.5pt}
\begin{tabular}{@{}lrrrrr@{}}
\toprule
 & \multicolumn{2}{c}{Simulated} & \multicolumn{3}{c}{Measured} \\
\cmidrule(lr){2-3}\cmidrule(l){4-6}
Trigger & dup.\,\% & rec.\,s & file & fBatch & rec.\,s \\
\midrule
\SimRows
\bottomrule
\end{tabular}
\end{table}

\section{Adaptive checkpoint selection, live}
\label{sec:acs}

\paragraph{The rule.}
ACS evaluates every 5\,s a risk score
$r = 0.4\,u + 0.3\,e + 0.3\,f$ from three signals in $[0,1]$: throughput
instability $u$ (coefficient of variation of input rate over the last ten
windows, saturating at 0.3; 0.2 during warm-up), error pressure $e$, and failure
recency $f = \max(0, 1 - t_{\text{since failure}}/120\,\text{s})$. If
$r > 0.45$ the interval shrinks by 30\%; otherwise the window counts as stable,
and after four stable windows the interval grows by 20\%, within 2--45\,s,
starting at 10\,s. In the simulation, $e$ counted exceptions per 1{,}000
records. A healthy live pipeline has none, so the live controller measures $e$
as the fraction of batches in the window whose run time exceeded the trigger
interval, a measurable sign of falling behind.

\paragraph{A structural flaw (v1).}
Implementing the rule exposed a property the simulation had hidden: no single
signal can cross the threshold. Instability alone contributes at most 0.4, and a
recent failure or full error pressure at most 0.38 (with warm-up instability).
And because every window with $r \le 0.45$ counts as stable, the interval
\emph{relaxes} during a burst unless a failure happened within the previous
100\,s. We therefore evaluate two variants, both fixed before any measured run:
\textbf{v1}, the rule exactly as specified, and \textbf{v2}, which tightens when
$r > 0.30$ (any one saturated signal), counts a window as stable only when
$r < 0.15$, and otherwise holds the interval.

\paragraph{Changing the interval of a running query.}
A running Structured Streaming query cannot change its trigger. Our controller
runs inside the Spark job and, to apply a new interval, waits until the query is
between batches, stops it, and starts it again on the same checkpoint in the
same Spark session, so no JVM restart is involved. Waiting for the gap between
batches matters: stopping a query interrupts a running batch, and for
foreachBatch a batch that had already written its output would be written again.
The controller persists its interval and a clean-stop flag next to the
checkpoint, so that after a crash the restarted job resumes the interval and
counts the crash for the recency signal.

\paragraph{Switching cost.}
Over \Switches{} switches in \ACSTrialsClean{} ACS trials, the time from stopping
the old query to the first batch of the new one had a median of \SwitchMed\,s
(\SwitchMin--\SwitchMax\,s; \SwitchOverOne{} of \SwitchN{} measured switches took
more than 1\,s), against about 10\,s for crash recovery. No ACS trial produced a
duplicate (\ACSDupTrials{} trials with duplicates). Restart-based adaptation is
therefore cheap and safe; whether it is \emph{useful} is the question of the next
two sections.

\section{Study 2: ACS against fixed intervals}
\label{sec:study2}

\paragraph{Design.}
Six strategies: fixed triggers of 1\,s (A) and 30\,s (B), ACS v1 and v2, and two
further fixed triggers set to ACS's mean intervals from the first run, 3.3\,s
and 5.8\,s (C). The C strategies are the fair comparison: they answer whether
\emph{adapting} helps, rather than merely choosing a value between 1\,s and
30\,s. The workload was bursty: every 15\,s the producer's rate was drawn from
0.5$\times$, 1$\times$ and 3$\times$ 300 records/s, from a seed equal to the trial
number. All trials used the foreachBatch sink and a driver kill at 15\,s plus a
phase spread over 30\,s, the same for every strategy. Trial $t$ therefore gave
every strategy the same input and the same kill time, so strategies can be
compared pair by pair.

\paragraph{Two runs.}
The first run (\ACSTrialsOne{} trials) ran A, B and both ACS variants
interleaved overnight, and the two C strategies the next morning, separately
from the others. The 5.8\,s C trials recovered more slowly (median 14.3\,s
against about 10\,s), and their Spark start-up was slower too, which no trigger
setting can cause, so the first run's
comparison with C is biased against C. The second run (\ACSTrialsTwo{} trials)
interleaved all six strategies trial by trial and is our main evidence. Six of
its trials were paused by laptop standby (Section~\ref{sec:stalls}); they are
kept in the data, excluded from the tables, and were rerun with the same trial
numbers, giving \ACSCleanTwo{} clean trials, 10 per strategy. Code commits:
\ACSCommits.

\begin{table*}[t]
\centering\small
\caption{Study 2, second (interleaved) run: per-trial medians, then median
(min--max) over 10 trials. Latency: producer send to output. Batch commits:
committed micro-batches per trial (36{,}000 records each). Mean interval:
time-weighted, ACS only.}
\label{tab:run2}
\resizebox{\textwidth}{!}{%
\begin{tabular}{@{}lrllllll@{}}
\toprule
Strategy & Trials & Latency p50 s & Latency p95 s & Batch commits & Mean interval s & Switches & Recovery s \\
\midrule
\RunTwoRows
\bottomrule
\end{tabular}}
\end{table*}

\begin{figure}[t]
\centering
\begin{tikzpicture}
\begin{axis}[width=\columnwidth, height=5.6cm, xmode=log, ymode=log,
  xlabel={Batch commits per trial (median)}, ylabel={Latency p50, s (median)},
  xmin=4, xmax=160, ymin=0.5, ymax=25, grid=major, grid style={gray!25},
  xtick={5,10,20,50,100}, xticklabels={5,10,20,50,100},
  ytick={0.5,1,2,5,10,20}, yticklabels={0.5,1,2,5,10,20},
  legend style={font=\small, at={(0.98,0.98)}, anchor=north east}]
\addplot[black, mark=*, mark size=1.6pt]
  table[x=commits, y=p50, restrict expr to domain={\thisrow{fixed}}{1:1}]
  {tables/frontier_run2.dat};
\addlegendentry{fixed interval}
\addplot[only marks, mark=triangle*, mark size=2.6pt, black!55]
  table[x=commits, y=p50, restrict expr to domain={\thisrow{fixed}}{0:0}]
  {tables/frontier_run2.dat};
\addlegendentry{ACS}
% generated by paper/make_tables.py from the committed trial records; do not edit
\node[font=\scriptsize, above] at (axis cs:82.0,0.843) {A (1 s)};
\node[font=\scriptsize, below left] at (axis cs:28.5,2.181) {C (3.3 s)};
\node[font=\scriptsize, above right] at (axis cs:35.0,1.986) {ACS v2};
\node[font=\scriptsize, below left] at (axis cs:18.0,3.465) {C (5.8 s)};
\node[font=\scriptsize, above right] at (axis cs:24.0,3.683) {ACS v1};
\node[font=\scriptsize, right] at (axis cs:6.0,13.977) {B (30 s)};

\end{axis}
\end{tikzpicture}
\caption{Latency against checkpoint overhead, second run. The line joins the
fixed intervals (B, C\,5.8\,s, C\,3.3\,s, A from left to right). ACS v2 lies on
the line; ACS v1 lies above it.}
\label{fig:frontier}
\end{figure}

\paragraph{Results.}
Table~\ref{tab:run2} and Figure~\ref{fig:frontier} show the second run. The
fixed intervals trace a clean trade-off: from A to B, median latency rises from
\PfiftyATwo\,s to \PfiftyBTwo\,s while batch commits fall from \CommitsATwo{} to
\CommitsBTwo. The question is where ACS lies relative to that curve.

\emph{ACS v1} (mean interval \IntervalAcsOneTwo\,s) had a median latency of
\PfiftyAcsOneTwo\,s with \CommitsAcsOneTwo{} commits; the fixed curve at the
same number of commits gives \CurveAcsOneTwo\,s. Against the fixed 5.8\,s
interval, it had a higher p95 latency (\PninetyfiveAcsOneTwo{} versus
\PninetyfiveCfiveTwo\,s) and more commits (\CommitsAcsOneTwo{} versus
\CommitsCfiveTwo). Pair by pair (Table~\ref{tab:paired}), v1 had the lower
median latency in \PairPfiftyAcsOneTwo{} pairs, the lower p95 in
\PairPninetyfiveAcsOneTwo, and fewer commits in \PairCommitsAcsOneTwo.

\emph{ACS v2} (mean interval \IntervalAcsTwoTwo\,s) had a median latency of
\PfiftyAcsTwoTwo\,s with \CommitsAcsTwoTwo{} commits, essentially on the fixed
curve (\CurveAcsTwoTwo\,s at that number of commits). Against the fixed 3.3\,s
interval it had a slightly lower median latency (\PfiftyAcsTwoTwo{} versus
\PfiftyCthreeTwo\,s; lower in \PairPfiftyAcsTwoTwo{} pairs), a higher p95
(\PninetyfiveAcsTwoTwo{} versus \PninetyfiveCthreeTwo\,s; lower in only
\PairPninetyfiveAcsTwoTwo) and more commits (\CommitsAcsTwoTwo{} versus
\CommitsCthreeTwo; fewer in only \PairCommitsAcsTwoTwo).

The first run (Table~\ref{tab:both}) shows the same ordering for A, B and the
two ACS variants. Its C strategies ran separately, and the 5.8\,s one had
higher latency than in the second run, which made ACS v1 look better against it
(\PairPfiftyAcsOneOne{} pairs with lower median latency); even so, every ACS
variant used more commits than the fixed interval at its own mean in both runs.

\begin{table}[t]
\centering\small
\caption{Paired comparison: in how many of the 10 pairs (same input, same kill
time) ACS had the lower value.}
\label{tab:paired}
\setlength{\tabcolsep}{3pt}
\begin{tabular}{@{}llccc@{}}
\toprule
Run & Comparison & p50 & p95 & Commits \\
\midrule
\PairRows
\bottomrule
\end{tabular}
\end{table}

\begin{table}[t]
\centering\small
\caption{Both runs, medians over clean trials. Run 1's C strategies ran
separately from the others (see text).}
\label{tab:both}
\setlength{\tabcolsep}{3pt}
\begin{tabular}{@{}lrrrrrr@{}}
\toprule
 & \multicolumn{3}{c}{Run 1} & \multicolumn{3}{c}{Run 2} \\
\cmidrule(lr){2-4}\cmidrule(l){5-7}
Strategy & p50 & p95 & comm. & p50 & p95 & comm. \\
\midrule
\BothRows
\bottomrule
\end{tabular}
\end{table}

\section{Why adaptation has nothing to gain}
\label{sec:mechanism}

\paragraph{Replay is rare.}
In the classic model, a failure loses the work done since the last checkpoint,
so a longer interval means more redone work. In micro-batch Spark, records that
arrived since the last batch were never processed; after a restart they are
read for the first time. Only a batch that was running when the driver died is
run again. Across both runs, that happened in \ReplayEvents{} of
\CleanTrials{} clean trials (Table~\ref{tab:replay}).

\paragraph{When it happens, it scales with the batch.}
The replayed amount is the size of the interrupted batch, so it grows with the
interval and the current rate: the largest, 15{,}750 records, came from a single 30\,s batch. Roughly, the chance of killing a batch in flight is its run
time divided by the interval, while its size grows with the interval, so the
expected replay depends much less on the interval than the classic model
suggests. Our ten trials per strategy are too few to estimate that expectation
precisely; we report the events individually.

\paragraph{Replay did not lengthen recovery, and rarely duplicated.}
Recovery in the replay trials was 9.0--14.0\,s, within the range of all other
trials; the 15{,}750-record replay recovered in 10.5\,s. Only one of the nine
replays produced duplicates (\DupRows{} rows): the interrupted batch had
already written its output, the at-least-once window from Study~1. In the
others the batch died before its write and was simply redone.

\begin{table}[t]
\centering\small
\caption{Every trial in which the kill interrupted a running batch (replay
$>0$), both runs. Dup.: duplicated rows in the output.}
\label{tab:replay}
\setlength{\tabcolsep}{4pt}
\begin{tabular}{@{}clrrrr@{}}
\toprule
Run & Strategy & Trial & Replayed & Recovery s & Dup. \\
\midrule
\ReplayRows
\bottomrule
\end{tabular}
\end{table}

\paragraph{Recovery does not depend on the interval.}
Median recovery was 10.0--10.5\,s for every strategy in the second run
(Table~\ref{tab:run2}). Splitting each recovery at the moment the restarted
query was running (Table~\ref{tab:split}), start-up took about 7\,s and the
first batch after it about 3.5\,s, for every strategy. Over both runs,
\RecFast{} of \RecAll{} recoveries took under 12\,s and \RecSlow{} took
12--17\,s; the slower ones had both slower start-up (median \SlowStartup\,s
against \FastStartup\,s) and a slower first batch (\SlowFirst{} against
\FastFirst\,s), which points to the machine's state at the time, not to the
strategy. One recovery took 25.6\,s, with an 18\,s start-up.

\begin{table}[t]
\centering\small
\caption{Recovery split, second run: Spark start-up after the restart, and the
first batch after start-up. Median (min--max), seconds.}
\label{tab:split}
\setlength{\tabcolsep}{4pt}
\begin{tabular}{@{}lrll@{}}
\toprule
Strategy & $n$ & Start-up & First batch \\
\midrule
\SplitRows
\bottomrule
\end{tabular}
\end{table}

\paragraph{What is left is latency against overhead.}
With replay rare and recovery independent of the interval, the interval governs
how long records wait for a batch (latency) and how many batches commit
(overhead). That trade-off depends on the workload's rate and the latency
target, both of which an operator knows in advance. A fixed interval chosen for
the target serves it directly. ACS adds switching (cheap, but not free), and its
decisions follow failure recency and throughput variance, which are signals for
the risk of losing work, the very cost shown above to be small here.

\section{Detecting paused trials}
\label{sec:stalls}

Six trials of the second run took minutes longer than any other, with p95
latencies up to 536\,s. Their outcomes were correct; the timings were not
Spark's. The producer is paced only by the wall clock, so the time it needs to
send all records follows from the workload schedule. We define a trial as
\emph{stalled} if its producer finished more than 10\,s later than its schedule
allows. Over \OverrunTotal{} trials, \OverrunUnderOne{} finished within 1\,s of
schedule, none between 1 and 10\,s, and \ACSStalled{} were 4 to 40 minutes late
(Table~\ref{tab:stalls}). The Windows power log showed that the laptop had
entered Modern Standby after an idle timeout during each of those periods,
although it was plugged in; an excerpt is in the repository. The test needs no
access to the operating system and catches a pause wherever it falls in a
trial; checking only the time to first output would have missed five of the six.

\begin{table}[t]
\centering\small
\caption{Stalled trials (producer more than 10\,s behind schedule), excluded
from the tables and rerun with the same trial numbers.}
\label{tab:stalls}
\begin{tabular}{@{}llrr@{}}
\toprule
Run id (start, UTC) & Strategy & Trial & Overrun s \\
\midrule
\StallRows
\bottomrule
\end{tabular}
\end{table}

\section{Threats to validity}
\label{sec:threats}

\begin{itemize}
\item \textbf{One machine, local mode.} Driver and executors share one JVM, so
executor-only failures, network partitions and multi-node scheduling are not
covered. On a cluster, batches take longer to schedule and the in-flight window
is larger, which could make replay more frequent; the mechanism in
Section~\ref{sec:mechanism} still applies, but its magnitude may differ.
\item \textbf{One workload, low rate.} 300 records/s with 0.5--3$\times$ bursts
and a trivial transformation. Heavier batches spend more time in flight, which
raises the chance that a kill interrupts one.
\item \textbf{Ten trials per strategy and run.} Enough to show the ordering
and its consistency across pairs and runs, not to estimate small differences
or rare-event rates. Replay events are reported individually for that reason.
\item \textbf{Machine state.} Recovery showed two modes tied to the machine's
state, and the first run's C strategies ran under different conditions. The
second run interleaves every strategy trial by trial; its six reruns ran
separately and recovered in the slower mode.
\item \textbf{One adaptive rule.} We tested ACS as specified and one revision.
A controller that targets latency directly, or tunes batch size as in
\cite{das2014adaptive}, is a different question; our results bound what an
interval controller driven by failure risk can gain here, not every possible
controller.
\item \textbf{One kind of corruption.} Overwriting part of the latest offsets
file; other corruptions (commit log, sink metadata, truncation) may behave
differently.
\end{itemize}

\section{Related work}

Discretized streams \cite{zaharia2013discretized} introduced micro-batch stream
processing with lineage-based recovery, and Structured Streaming
\cite{armbrust2018structured} its declarative successor with exactly-once sinks.
Das et al.\ \cite{das2014adaptive} adapt the batch interval of Spark Streaming
to keep the system stable as load changes, a latency and stability goal rather
than a fault-tolerance one, consistent with our finding that the interval is a
latency knob. Drizzle \cite{venkataraman2017drizzle} decouples the processing
interval from coordination to cut micro-batch overhead. Flink takes consistent
snapshots with barriers derived from Chandy--Lamport
\cite{chandy1985snapshots,carbone2015lightweight,carbone2017state}, where the
checkpoint interval and the processing granularity are separate and the classic
trade-off applies more directly. Optimal checkpoint intervals under failure
were derived by Young \cite{young1974first} and Daly \cite{daly2006higher}, and
rollback recovery is surveyed by Elnozahy et al.\ \cite{elnozahy2002survey}.
The same measurement approach, a fault-injection harness with committed raw
logs, was applied to LLM serving in \cite{medipally2026llm}.

\section{Conclusion}

In Spark Structured Streaming on one machine, a driver crash rarely forced a
batch to be redone (\ReplayEvents{} of \CleanTrials{} trials), because records
that arrive between batches are read after the restart, not replayed. Recovery
took about 10\,s regardless of the trigger interval, mostly in restart. The
interval therefore traded latency against checkpoint overhead. Changing it at
runtime was cheap and safe, but on paired workloads and failures, a fixed
interval at ACS's own mean matched or beat it, and the original ACS rule could
not react to any single signal. For micro-batch Spark, our evidence supports
choosing a fixed trigger for the latency target over adapting it to failure
risk. Next steps: a multi-node cluster with executor failures and heavier
batches, where in-flight batches are larger; other corruption modes; and
controllers that target latency directly.

\paragraph{Artifact.}
Code, every trial record and log, and the script that generates every table,
figure and in-text number of this paper are at
\url{https://github.com/rmedipallycic/spark-streaming-fault-tolerance}
(MIT licence). Excluded trials are kept with their reasons.

\section*{Acknowledgments}
The harness code, analysis scripts, and much of the text of this paper were
developed with substantial help from an AI assistant (Claude, Anthropic). The
author directed the work, ran all experiments, and is responsible for every
reported result; each number is generated by \texttt{paper/make\_tables.py}
from the trial records committed to the repository.

\clearpage
\bibliographystyle{plain}
\bibliography{refs}
\end{document}
